\begin{abstract}
Die automatisierte Klassifikation von Fernerkundungsbildern ist ein zentrales Problem in der Geoinformatik und findet Anwendung in Bereichen wie der Landnutzungskartierung, Stadtplanung und Umweltüberwachung.
In dieser Arbeit werden zehn Convolutional-Neural-Network-Architekturen systematisch entworfen, trainiert und verglichen, um Szenen des PatternNet-Datensatzes in 38~Landnutzungsklassen einzuordnen.
Ausgehend von einem schlanken Basisnetzwerk werden schrittweise einzelne Komponenten, wie zusätzliche Convolution-Layer, aggressiveres Pooling, Dropout und Batch-Normalisierung, isoliert variiert, bevor die wirksamsten Techniken in kombinierte Architekturen integriert werden.
Alle Netze werden unter identischen Bedingungen trainiert (Adam-Optimizer, 15~Epochen, Batchgröße~64) und anhand der Klassifikationsgenauigkeit auf einem unabhängigen Testdatensatz bewertet.
Das leistungsfähigste Modell erreicht eine Test-Accuracy von \textbf{84\,2\%}.
Die Ergebnisse zeigen, dass Batch-Normalisierung den größten Einzeleffekt auf die Klassifikationsleistung hat, während zusätzliche Faltungsschichten oder tiefere Klassifikationsköpfe allein nur marginale Verbesserungen erzielen, und verdeutlichen damit die Bedeutung einer gezielten Regularisierung gegenüber einer reinen Erhöhung der Netztiefe.
\end{abstract}